\newpage
\section{极坐标}
\label{sec:quadratic-polar-coordinates}

直角坐标用两个方向上的分量确定点的位置，也可以用到定点的距离和方向来确定位置。这两种表示描述同一个平面，却适合不同的问题。本节先建立极坐标，再按研究对象选择坐标原点，用距离与角度重新表达二次曲线、极线和焦弦的关系。

\subsection{极坐标系与坐标转换}
\label{subsec:quadratic-coordinates}

\begin{definition}
在平面上选定一点\(O\)、从\(O\)出发的一条射线和长度单位，分别称\(O\)为{\heiti{}极点}，该射线为{\heiti{}极轴}。对\(P\ne O\)，记\(r=|OP|\)，从极轴逆时针转到射线\(OP\)的角为\(\theta\)，则\((r,\theta)\)称为\(P\)的{\heiti{}极坐标}。其中\(r\)称为极径，\(\theta\)称为极角。
\end{definition}

本节通常取\(r\ge0\)，角以弧度为单位。角度相差\(2k\pi\)（\(k\in\mathbb Z\)）给出同一方向，因此同一点的极坐标不唯一。极点对应\(r=0\)，此时极角任意。若需使非极点的表示唯一，可以将极角限制在\([0,2\pi)\)等半开区间内。

以极点为直角坐标原点、极轴为\(x\)轴正方向，建立同一长度单位的直角坐标系。由三角函数的定义，两个坐标之间有
\begin{equation}\label{eq:quadratic-polar-cartesian-conversion}
 x=r\cos\theta,\qquad y=r\sin\theta;
 \qquad r^2=x^2+y^2.
\end{equation}
由\((x,y)\)求极角时，还须用坐标的正负确定所在象限，不能只由\(\tan\theta=y/x\)确定方向。

\begin{figure}[!ht]
 \centering
 \includegraphics[width=.60\linewidth]{conic-polar-coordinates}
 \caption{用距离与方向确定点的位置}
 \label{fig:quadratic-polar-coordinates}
\end{figure}

若极点在原直角坐标的\((h,k)\)，极轴与原\(x\)轴正向成角\(\alpha\)，则
\begin{equation}\label{eq:quadratic-shifted-polar-conversion}
 x=h+r\cos(\alpha+\theta),\qquad
 y=k+r\sin(\alpha+\theta).
\end{equation}
这相当于先平移、旋转直角坐标，再使用式\eqref{eq:quadratic-polar-cartesian-conversion}。选定极点和极轴是建系的一部分，可以根据曲线的中心、焦点或问题中的定点来决定。

两点\(P_1(r_1,\theta_1),P_2(r_2,\theta_2)\)之间的距离，由余弦定理或坐标计算得到
\begin{equation}\label{eq:quadratic-polar-distance}
 |P_1P_2|^2=r_1^2+r_2^2-2r_1r_2\cos(\theta_1-\theta_2).
\end{equation}
特别地，夹角为直角时，右边成为\(r_1^2+r_2^2\)。

\paragraph{1°曲线方程的转换} \mbox{}

曲线上的点在两种坐标下相同，将式\eqref{eq:quadratic-polar-cartesian-conversion}代入直角坐标方程，就得到极坐标方程。反过来，利用\(r\cos\theta=x\)、\(r\sin\theta=y\)及\(r^2=x^2+y^2\)，也可以还原直角坐标方程。

例如，圆心在极点、半径为\(R\)的圆表示为\(r=R\)。圆\((x-a)^2+y^2=a^2\)（\(a>0\)）则给出
\begin{equation*}
 r^2-2ar\cos\theta=0.
\end{equation*}
非极点部分满足\(r=2a\cos\theta\)，取\(-\pi/2\le\theta\le\pi/2\)可表示整个圆，端点角度给出极点。这说明约去\(r\)后，应检查极点是否仍由新方程表示。

对不过极点的直线，取其单位法向量方向角为\(\beta\)，方程可写成
\begin{equation}\label{eq:quadratic-polar-line-normal}
 x\cos\beta+y\sin\beta=d
 \quad\Longleftrightarrow\quad
 r\cos(\theta-\beta)=d.
\end{equation}
选择法向量使\(d>0\)时，\(d\)就是极点到直线的距离。过极点、方向角为\(\alpha\)的整条直线，由相反的两个方向\(\theta=\alpha\)、\(\theta=\alpha+\pi\)组成。

\par\addvspace{6pt}\noindent\begin{minipage}{\linewidth}
\begin{example}
分别将直线\(x+y=2\)、圆\(x^2+y^2-4x=0\)、曲线\(x^2-y^2=4\)写成极坐标方程，并说明各自允许的极角范围。再将\(r\cos(\theta-\pi/3)=2\)还原为直角坐标方程。
\end{example}
\end{minipage}\par\addvspace{6pt}

\paragraph{2°沿直线的有向距离} \mbox{}

研究一条过极点的直线时，取单位方向向量\(d=(\cos\theta,\sin\theta)\)，将线上点写成
\begin{equation}\label{eq:quadratic-polar-signed-radius}
 X=O+td,\qquad t\in\mathbb R.
\end{equation}
\begin{samepage}
\(t\)是沿该方向的有向距离：\(t>0\)对应极角\(\theta\)，\(t<0\)对应极角\(\theta+\pi\)，实际长度为\(|t|\)。这也可以理解为允许负极径，并约定\((t,\theta)\)与\((-t,\theta+\pi)\)表示同一点。下面研究整条割线时使用\(t\)，研究射线上的距离时仍使用\(r\ge0\)。
\par\end{samepage}

\par\addvspace{6pt}\noindent\begin{minipage}{\linewidth}
\begin{example}
以\((1,-1)\)为极点、平行于原\(x\)轴正向的射线为极轴，将\((x-1)^2+2(y+1)^2=4\)写成极坐标方程。求曲线上极角相差\(\pi/2\)的两点到极点距离的平方倒数之和，判断它是否为定值。
\end{example}
\end{minipage}\par\addvspace{6pt}

\begin{exercise}
\mbox{}\par\nobreak
\begin{samepage}
（1）写出\((-\sqrt3,1)\)、\((0,-2)\)在\(0\le\theta<2\pi\)内的极坐标，再求极坐标为\((2,5\pi/6)\)、\((3,-\pi/2)\)的点的直角坐标。
\par\end{samepage}
\begin{samepage}
（2）将\(r=4\sin\theta\)、\(r^2\cos2\theta=2\)、\(r(2\cos\theta-\sin\theta)=3\)化为直角坐标方程，说明各自表示的曲线。
\par\end{samepage}
\begin{samepage}
（3）以中心为极点，写出椭圆\(x^2/9+y^2/4=1\)的极坐标方程，求极径范围及取等方向。
\par\end{samepage}
\begin{samepage}
（4）两点的极坐标为\((2,\pi/6)\)、\((3,5\pi/6)\)。求两点距离和它们与极点组成的三角形面积。
\par\end{samepage}
\begin{samepage}
（5）将圆\((x-2)^2+(y+1)^2=9\)分别写成以原点、以圆心为极点的极坐标方程，比较两种表示。极轴均取平行于原\(x\)轴正向的射线。
\par\end{samepage}
\begin{samepage}
（6）（据2023年全国乙卷理科第22题）以原点为极点、\(x\) 轴正向为极轴。曲线 \(C_1:r=2\sin\theta\)（\(\pi/4\le\theta\le\pi/2\)），\(C_2:(x,y)=(2\cos\alpha,2\sin\alpha)\)（\(\pi/2<\alpha<\pi\)）。写出 \(C_1\) 的直角坐标方程；直线 \(y=x+m\) 与两曲线均不相交时，\(m\) 取何值？
\par\end{samepage}

\begin{samepage}
（7）圆的极坐标方程为\(r^2-4r\cos\theta+3=0\)。求有实交点的射线方向，判断哪些射线与圆有两个不同交点，并说明这条曲线能否在全部这些方向上写成单值的\(r=f(\theta)\)。
\par\end{samepage}
\begin{samepage}
（8）以椭圆\(x^2/9+y^2/4=1\)的中心为极点，将极轴取为与\(x\)轴正向成\(\pi/4\)角的射线。写出极坐标方程，求极角为\(0\)、\(\pi/2\)的两点的直角坐标。
\par\end{samepage}
\begin{samepage}
（9）在圆\(x^2+y^2=4\)上取\(A=(2,0)\)为极点，极轴指向圆心。写出圆的极坐标方程，并用它证明：若圆上两点\(B,C\)满足\(\angle BAC=\pi/3\)，则\(|BC|\)为定值。
\par\end{samepage}

\begin{samepage}
（10）椭圆\(x^2/a^2+y^2/b^2=1\)（\(a>b>0\)）的中心为\(O\)。曲线上两点\(A,B\)满足\(OA\perp OB\)，用中心极坐标求\(\triangle OAB\)面积的最大值、最小值及相应方向，并与共轭半径的面积关系比较。
\par\end{samepage}
\begin{samepage}
（11）点 \(P\) 在直线 \(x+2y=4\) 上移动，射线 \(OP\) 交椭圆 \(x^2/4+y^2=1\) 于 \(R\)。在该射线上取 \(Q\) 使 \(|OP||OQ|=|OR|^2\)，求 \(Q\) 的轨迹。
\par\end{samepage}
\end{exercise}

\newpage
\subsection{极线与径向方程}
\label{subsec:quadratic-polar-line-and-radial-equation}

极坐标系的极点可以任意选取；点关于二次曲线的极线，则由曲线和该点决定。现在选定一点\(O\)作为极坐标原点，研究它的极线怎样进入沿各个方向的交点方程。

平移后以\(O\)为直角坐标原点，将所研究的非退化二次曲线写为
\begin{equation}\label{eq:quadratic-polar-origin-equation}
 Ax^2+Bxy+Cy^2+Dx+Ey+F=0.
\end{equation}
沿方向\(\theta\)代入\(x=t\cos\theta,y=t\sin\theta\)，得到
\begin{equation}\label{eq:quadratic-radial-scalar-equation}
 a(\theta)t^2+b(\theta)t+F=0,
\end{equation}
其中
\begin{equation*}
 \begin{aligned}
 a(\theta)&=A\cos^2\theta+B\cos\theta\sin\theta+C\sin^2\theta,\\
 b(\theta)&=D\cos\theta+E\sin\theta.
 \end{aligned}
\end{equation*}
这就是该方向上的{\heiti{}径向方程}。取\(t\ge0\)时得到相应射线上的曲线点；允许\(t\)取负值时，同时研究整条直线的交点。

根据4.10的极线公式，原点\(O=(0,0)\)的极线为
\begin{equation}\label{eq:quadratic-radial-polar-line}
 Dx+Ey+2F=0.
\end{equation}
可见径向方程的一次项由极线决定，而二次项还包含曲线在各个方向上的信息。只知道极线，通常还不足以确定整条曲线。

\paragraph{1°交点的倒数关系} \mbox{}

设\(O\)不在曲线上，一条过\(O\)的直线与曲线交于两个不同的有限点\(U,V\)，有向距离为\(t_1,t_2\)。由根与系数的关系，
\begin{equation}\label{eq:quadratic-radial-reciprocal-sum}
 \frac1{t_1}+\frac1{t_2}=-\frac{b(\theta)}F.
\end{equation}
若这条直线与\(O\)的极线交于\(H\)，设其有向距离为\(h\)，由式\eqref{eq:quadratic-radial-polar-line}得\(b(\theta)h+2F=0\)。所以
\begin{equation}\label{eq:quadratic-radial-harmonic-mean}
 \boxed{\ \frac1{t_1}+\frac1{t_2}=\frac2h\ }.
\end{equation}
这正是4.10中割线与极线的距离关系。现在看出，它是径向二次方程的根与系数关系，与坐标原点取在何处无关。

若\(b(\theta)=0\)，两根互为相反数；当\(O\)有极线时，该方向与极线平行。

中心处\(D=E=0\)，各方向的交点均关于中心成对。

有向距离不能直接替换为长度。例如，\(U,V\)在\(O\)两侧时，\(t_1,t_2\)异号，上式对应两个长度倒数之差；在同侧时则对应倒数之和。

\par\addvspace{6pt}\noindent\begin{minipage}{\linewidth}
\begin{example}
对椭圆\(x^2/9+y^2/4=1\)，取\(O=(1,0)\)为极点，极轴平行于\(x\)轴正向。写出径向方程及\(O\)的极线，分别沿水平、竖直方向计算交点，核对式\eqref{eq:quadratic-radial-reciprocal-sum}。水平割线上的有向距离与实际长度有什么不同？
\end{example}
\end{minipage}\par\addvspace{6pt}

\paragraph{2°由交点研究方向} \mbox{}

当\(a(\theta)\ne0\)时，径向方程的判别式为\(\Delta(\theta)=b(\theta)^2-4a(\theta)F\)，其正、零、负分别表示该直线有两个不同实交点、在一个点相切、没有实交点。若只研究某条射线，还要看所得根是否为正。若\(a(\theta)=0\)，则直接研究剩余方程的有限交点。

这与4.8的联立方法相同，而方向角也包括了竖直方向。计算时可按任务选取根、根的和与积，或判别式。

\par\addvspace{6pt}\noindent\begin{minipage}{\linewidth}
\begin{example}
对圆\(x^2+y^2=1\)，以\(O=(2,0)\)为极点，极轴取从\(O\)指向圆心的方向。求径向方程、能与圆相交的射线方向，以及两条切线的夹角。再求每条实割线的两交点到\(O\)的距离之积，并与圆幂比较。
\end{example}
\end{minipage}\par\addvspace{6pt}

\begin{exercise}
\mbox{}\par\nobreak
\begin{samepage}
（1）曲线为\(x^2+2y^2=4\)，极点取\(O=(1,1)\)，极轴平行于\(x\)轴正向。写出径向方程及\(O\)的极线，并求方向角\(\pi/4\)的直线与极线的交点。
\par\end{samepage}
\begin{samepage}
（2）对抛物线\(y^2=4x\)，取\(O=(-1,0)\)。求从\(O\)出发的切线方向，并写出两条切线的方程。
\par\end{samepage}
\begin{samepage}
（3）在椭圆\(x^2/9+y^2/4=1\)中，\(O=(1,0)\)。证明过\(O\)的各条弦，两个端点到\(O\)的距离之积为\(8/(\cos^2\theta+\tfrac94\sin^2\theta)\)，并求其范围。这里\(\theta\)是弦的方向角。
\par\end{samepage}
\begin{samepage}
（4）过双曲线\(x^2-y^2=1\)上点\(O=(1,0)\)作直线。直接从径向方程分离恒有的根\(t=0\)，求另一交点的有向距离，说明相切方向与平行于渐近线的方向在计算中有何不同。
\par\end{samepage}
\begin{samepage}
（5）设\(O\)不在非退化二次曲线上，沿两条不同方向的实割线分别取交点有向距离的倒数和。若这两个和都为零，证明\(O\)是曲线的中心。
\par\end{samepage}
\begin{samepage}
（6）对椭圆\(x^2/a^2+y^2/b^2=1\)，以中心为极点，求方向角为\(\theta\)的半径长\(r(\theta)\)。证明\(r(\theta)^{-2}+r(\theta+\pi/2)^{-2}\)为定值，并解释相应两条半径互相垂直。
\par\end{samepage}
\end{exercise}

\newpage
\subsection{焦点极坐标方程}
\label{subsec:quadratic-focal-polar-equation}

极点若取在焦点，点到焦点的距离就成为极径，焦准线关系因而可以直接用于写方程。这一选择与上小节相连：4.10已经证明，非圆圆锥曲线的焦点与相应准线互为极点、极线。

\paragraph{1°由焦准线关系建立方程} \mbox{}

取焦点\(S\)为极点，从\(S\)垂直指向对应准线的射线为极轴。设焦点到准线的距离为\(d>0\)，离心率为\(e>0\)。在相应的直角坐标中，\(S=(0,0)\)，准线为\(x=d\)。曲线上点\(P(r,\theta)\)满足
\begin{equation}\label{eq:quadratic-focal-polar-absolute}
 r=e|d-r\cos\theta|.
\end{equation}
记\(\ell=ed\)。当\(P\)与焦点在准线同侧时，上式成为\(r=e(d-r\cos\theta)\)，故
\begin{equation}\label{eq:quadratic-focal-polar-equation}
 \boxed{\ r=\frac{\ell}{1+e\cos\theta}\ },
 \qquad 1+e\cos\theta>0.
\end{equation}
将\(\theta=\pi/2\)或\(3\pi/2\)代入，得\(r=\ell\)。对应两点组成通径，所以\(\ell\)是通径长的一半，称为{\heiti{}半通径}。由4.5中的通径公式，
\begin{equation}\label{eq:quadratic-semilatus-rectum}
 \ell=\frac{b^2}{a}\quad\text{（椭圆、双曲线）},
 \qquad \ell=p\quad\text{（抛物线\(y^2=2px\)）}.
\end{equation}

\begin{figure}[!ht]
 \centering
 \includegraphics[width=.65\linewidth]{conic-focal-polar}
 \caption{以焦点为极点、朝向准线的方向为极轴}
 \label{fig:quadratic-focal-polar}
\end{figure}

当\(0<e<1\)时，分母对所有角度均为正，式\eqref{eq:quadratic-focal-polar-equation}表示整条椭圆；当\(e=1\)时，取\(\theta\ne\pi+2k\pi\)，表示整条抛物线。在后一情形中，\(\theta=\pi\)方向的射线没有曲线点。

当\(e>1\)时，还可能有点在准线另一侧，此时
\begin{equation}\label{eq:quadratic-hyperbola-second-polar-radius}
 r=\frac{\ell}{e\cos\theta-1},
 \qquad e\cos\theta>1.
\end{equation}
因此双曲线的完整表示需合用式\eqref{eq:quadratic-focal-polar-equation}与式\eqref{eq:quadratic-hyperbola-second-polar-radius}。某些方向的射线会与双曲线的两支各交一次，极径不再是极角的单值函数。

若将极轴反向，式中的\(\cos\theta\)随之换号；若将极轴旋转，则相应改变角度。常见的\(1-e\cos\theta\)、\(1+e\sin\theta\)等形式，由同一焦准线关系在不同建系下得到。

将式\eqref{eq:quadratic-focal-polar-absolute}平方，再使用坐标转换，得到
\begin{equation}\label{eq:quadratic-focal-polar-cartesian}
 (1-e^2)x^2+y^2+2e\ell x-\ell^2=0.
\end{equation}
反过来，这个方程等价于\(r^2=e^2(d-x)^2\)，因两边的距离非负，仍等价于焦准线关系。这使离心率分类、直角坐标方程与极坐标表示接在了一起。

\par\addvspace{6pt}\noindent\begin{minipage}{\linewidth}
\begin{example}
曲线的极坐标方程为\(r=6/(2+\cos\theta)\)。求焦点、对应准线、离心率和通径长，再求直角坐标标准方程以及另一焦点的位置。
\end{example}
\end{minipage}\par\addvspace{6pt}

\par\addvspace{6pt}\noindent\begin{minipage}{\linewidth}
\begin{example}
焦点为原点、准线为\(x=3\)、离心率为\(2\)的双曲线，以\(x\)轴正向为极轴。写出两种正极径表达，求\(\theta=0\)、\(\pi/2\)、\(\pi\)时射线上的全部交点。将焦准线关系还原为标准方程，核对两支的位置。
\end{example}
\end{minipage}\par\addvspace{6pt}

\paragraph{2°焦弦与方向} \mbox{}

对椭圆和抛物线，设焦弦\(UV\)的两端极角为\(\theta,\theta+\pi\)，相应焦半径为\(r_1,r_2\)，则
\begin{equation*}
 r_1=\frac{\ell}{1+e\cos\theta},\qquad
 r_2=\frac{\ell}{1-e\cos\theta}.
\end{equation*}
因此
\begin{align}
 \frac1{r_1}+\frac1{r_2}&=\frac2\ell,
 \label{eq:quadratic-focal-polar-reciprocals}\\
 |UV|=r_1+r_2&=\frac{2\ell}{1-e^2\cos^2\theta}.
 \label{eq:quadratic-focal-polar-chord-length}
\end{align}
抛物线中焦弦不能平行于轴，式中的分母自然排除了这一方向。椭圆的焦弦长介于\(2\ell\)与\(2\ell/(1-e^2)=2a\)之间，最短者为通径，最长者为长轴。抛物线的焦弦长为\(2p/\sin^2\theta\)，最小值仍在通径处取得。

用式\eqref{eq:quadratic-focal-polar-reciprocals}回看例4.39：知道焦点分弦的比，能否直接把两个焦半径写成半通径的倍数？

这里的倒数和与上小节的有向倒数和不同。两端极角相差\(\pi\)，实际长度都为正；若在整条直线上采用同一个正方向，另一端的有向距离应取负值。

双曲线仍可按焦弦方向分两种情况。若\(|e\cos\theta|<1\)，两个端点在焦点两侧，上述两式仍成立，且它们位于同一支上。若\(|e\cos\theta|>1\)，可反向选取\(\theta\)使\(e\cos\theta>1\)，两个交点在同一射线上，其距离为
\begin{equation*}
 r_1=\frac{\ell}{e\cos\theta+1},\qquad
 r_2=\frac{\ell}{e\cos\theta-1}.
\end{equation*}
此时\(1/r_1-1/r_2=2/\ell\)，而弦长为
\begin{equation}\label{eq:quadratic-hyperbola-focal-polar-chord}
 |UV|=r_2-r_1=\frac{2\ell}{e^2\cos^2\theta-1}.
\end{equation}
当\(|e\cos\theta|=1\)时，该直线平行于一条渐近线，只有一个有限交点。两种焦弦的长度关系不同，原因正是焦点在线段内部或外部的位置不同。

\par\addvspace{6pt}\noindent\begin{minipage}{\linewidth}
\begin{example}
椭圆\(x^2/9+y^2/5=1\)的右焦点为\(S\)。过\(S\)的弦\(UV\)满足\(|UV|=5\)，求弦与长轴夹角的余弦绝对值，以及两段焦半径的长度。分别考虑如何使用焦点极坐标和直角坐标处理。
\end{example}
\end{minipage}\par\addvspace{6pt}

\paragraph{3°垂直方向与距离最值} \mbox{}

若抛物线\(y^2=2px\)的一条焦弦与轴的夹角为\(\theta\)，另一条与之垂直，则两弦长度为
\begin{equation*}
 L_1=\frac{2p}{\sin^2\theta},\qquad
 L_2=\frac{2p}{\cos^2\theta}.
\end{equation*}
所以\(1/L_1+1/L_2=1/(2p)\)，并可进一步研究长度和、长度积的最小值。角度把垂直条件变成了\(\pi/2\)的差，再由三角恒等式进入计算。

极坐标也适用于过焦点的三角形。若两条焦半径夹角为\(\varphi\)，面积为\(r_1r_2|\sin\varphi|/2\)；若需要第三边，则使用式\eqref{eq:quadratic-polar-distance}。是否简便，要看距离、角度条件在所选坐标中能否直接表达。

\par\addvspace{6pt}\noindent\begin{minipage}{\linewidth}
\begin{example}
抛物线\(y^2=4x\)的焦点为\(S\)，过\(S\)作两条互相垂直的焦弦，长度分别为\(L_1,L_2\)。求\(L_1+L_2\)、\(L_1L_2\)的最小值，并说明等号时两条弦的方向。
\end{example}
\end{minipage}\par\addvspace{6pt}


\par\addvspace{6pt}\noindent\begin{minipage}{\linewidth}
\begin{example}\label{ex:quadratic-focal-projection-minimum}
椭圆\(x^2/3+y^2=1\)的左焦点为\(S\)。过\(S\)且斜率为正的直线交椭圆于\(A,B\)，过两端分别作弦的垂线，与\(x\)轴交于\(M,N\)。用焦点极坐标求\(|MN|\)的最小值，并说明等号时焦弦的方向。
\end{example}
\end{minipage}\par\addvspace{6pt}

\newpage
\paragraph{4°弦的方程与角度关系} \mbox{}

极坐标也可以表示不经过极点的弦。以椭圆\(r=\ell/(1+e\cos\theta)\)为例，设两个不同端点的极角为\(\alpha,\beta\)，记
\begin{equation*}
 m=\frac{\alpha+\beta}{2},\qquad
 \delta=\frac{\alpha-\beta}{2}.
\end{equation*}
对任一端点，\(\cos(\theta-m)=\cos\delta\)，且\(r+er\cos\theta=\ell\)，所以两端都满足
\begin{equation}\label{eq:quadratic-focal-polar-chord}
 r\bigl[\cos(\theta-m)+e\cos\delta\cos\theta\bigr]
 =\ell\cos\delta.
\end{equation}
左边使用\(r\cos\theta=x,r\sin\theta=y\)展开，成为
\begin{equation*}
 (\cos m+e\cos\delta)x+\sin m\,y=\ell\cos\delta.
\end{equation*}
这是经过两个端点的直线，因而就是弦所在直线。\(\cos\delta=0\)时两端方向相反，式子相应表示过焦点的弦。推导的关键是先用焦点方程替换极径，再将两端的角度关系写入一次方程。

\par\addvspace{6pt}\noindent\begin{minipage}{\linewidth}
\begin{example}\label{ex:quadratic-focal-cross-chord-angles}
椭圆的两条焦弦\(AB,CD\)交于同一焦点\(S\)，四端点互异。设\(U=AC\cap BD\)、\(V=AD\cap BC\)均为有限点。以\(S\)为极点，用弦的极坐标方程证明\(SU,SV\)分别沿两条焦弦所在直线的两条角平分线，因而\(SU\perp SV\)。再求\(U,V\)所在定直线，并与极点极线方法比较。
\end{example}
\end{minipage}\par\addvspace{6pt}

以\(SA,SC\)的方向角为\(\alpha,\beta\)，则\(SB,SD\)分别增加\(\pi\)。两组弦方程中的和角、差角由此相互联系；相加可以消去方向项，相减可以消去常数项。这个配置中，准线可以由坐标计算或极线关系找到，角平分关系则进一步使用了方向角的信息。

\begingroup
\setlength{\parskip}{4pt}
\begin{exercise}
\mbox{}\par\nobreak
\begin{samepage}
（1）求\(r=4/(2+\sin\theta)\)的离心率、对应准线与直角坐标方程，求曲线上点到极点距离的范围。
\par\end{samepage}
\begin{samepage}
（2）对椭圆\(x^2/25+y^2/9=1\)，分别以左、右焦点为极点、原\(x\)轴正向为极轴，写出极坐标方程。说明两个方程中余弦项的符号为何不同。
\par\end{samepage}
\begin{samepage}
（3）曲线满足\(r=3/(1+\cos\theta)\)。求直角坐标方程、顶点与焦点；再求过极点、与极轴成\(\pi/3\)角的焦弦长。
\par\end{samepage}
\begin{samepage}
（4）椭圆的离心率为\(1/2\)，通径长为\(6\)。求长、短半轴长。若某焦弦被焦点分成长度比为\(2:1\)的两段，求两段长度及弦与长轴的夹角。
\par\end{samepage}
\begin{samepage}
（5）双曲线\(x^2/4-y^2/5=1\)的左焦点为\(S\)。求过\(S\)且与实轴成\(\pi/3\)角的焦弦长，判断两端位于同一支还是不同支；再对实轴方向作同样判断。
\par\end{samepage}
\begin{samepage}
（6）抛物线\(y^2=2px\)的一条焦弦两端为\(U,V\)，证明\(|UV|=2|SU||SV|/p\)。若\(|SU|:|SV|=3:1\)，求弦长及其与轴夹角的余弦绝对值。
\par\end{samepage}
\begin{samepage}
（7）在椭圆上取三条过同一焦点的弦，方向角依次为\(\theta,\theta+\pi/3,\theta+2\pi/3\)，长度为\(L_1,L_2,L_3\)。证明\(1/L_1+1/L_2+1/L_3\)为定值，并用离心率和半通径表示这个定值。
\par\end{samepage}
\begin{samepage}
（8）（据2023年全国甲卷理科第20题）设\(F\)为抛物线\(y^2=2px\)（\(p>0\)）的焦点，曲线上两点\(M,N\)满足\(\overrightarrow{FM}\cdot\overrightarrow{FN}=0\)。用\(p\)表示\(\triangle MFN\)面积的最小值，并指出取等时两条射线的方向。若直线\(x-2y+1=0\)被抛物线截得的弦长为\(4\sqrt{15}\)，这个最小值是多少？
\par\end{samepage}

\begin{samepage}
（9）椭圆\(x^2/16+y^2/7=1\)的右焦点为\(S\)。曲线上两点\(P,Q\)满足\(SP\perp SQ\)，求\(1/|SP|+1/|SQ|\)的取值范围，并说明等号时两条射线的方向。
\par\end{samepage}
\begin{samepage}
（10）椭圆的长、短半轴长分别为\(a,b\)，焦点为\(S_1,S_2\)。用焦点极坐标求曲线上点\(P\)所对应的\(|PS_1||PS_2|\)的最大值与最小值，以及相应点的位置。
\par\end{samepage}
\begin{samepage}
（11）椭圆的离心率为\(e\)，半通径为\(\ell\)。过同一焦点作两条互相垂直的焦弦，长为\(L_1,L_2\)。证明\(1/L_1+1/L_2\)为定值，进一步求\(L_1L_2\)的取值范围。
\par\end{samepage}
\begin{samepage}
（12）抛物线\(y^2=2px\)的顶点为\(V\)，焦弦\(AB\)的长度为\(8p\)。求焦弦与轴的夹角及\(\triangle VAB\)的面积。
\par\end{samepage}
\begin{samepage}
（13）对双曲线\(x^2/a^2-y^2/b^2=1\)，分别研究两端在同一支、两端在不同支的焦弦。证明两类焦弦的最短长度分别为\(2b^2/a\)、\(2a\)，并指出相应方向。
\par\end{samepage}

\begin{samepage}
（14）非圆圆锥曲线的离心率为\(e\)，焦点为\(S\)。一条焦弦\(AB\)的中垂线与焦点所在的对称轴交于唯一点\(D\)。证明\(|SD|=e|AB|/2\)。用焦点极坐标分别处理\(S\)在线段\(AB\)内部和外部的情形。
\par\end{samepage}
\begin{samepage}
（15）双曲线的离心率为\(e\)，半通径为\(\ell\)。过同一焦点作两条互相垂直的实焦弦，长为\(L_1,L_2\)。令\(\varepsilon_i=1\)或\(-1\)，分别对应两端在同一支或不同支，证明\(\varepsilon_1/L_1+\varepsilon_2/L_2=(2-e^2)/(2\ell)\)。再取\(e=2,\ell=6\)，比较方向角为\(0\)与\(\pi/4\)时两条垂直焦弦的实际长度倒数和。
\par\end{samepage}
\begin{samepage}
（16）（2022年新高考I卷第16题）椭圆的上顶点为\(A\)，左、右焦点分别为\(F_1,F_2\)，离心率为\(1/2\)。过\(F_1\)且垂直于\(AF_2\)的直线与椭圆交于\(D,E\)。已知\(|DE|=6\)，求\(\triangle ADE\)的周长。
\par\end{samepage}
\end{exercise}
\endgroup

\newpage
\subsection{二次型与径向方程*}
\label{subsec:quadratic-quadratic-forms-and-radial-equation}

前面的径向方程由坐标代入得到。现在将原点、方向与二次曲线分开，用二次型说明这些系数各自表示什么。

设在原直角坐标中，曲线方程为
\begin{equation*}
 \Phi(x)=x^THx+L^Tx+F=0,\qquad H=H^T.
\end{equation*}
选点\(P\)为极坐标原点，记\(d\)为某一方向的单位向量。展开\(\Phi(P+td)\)，得到
\begin{equation}\label{eq:quadratic-radial-matrix-equation}
 \Phi(P+td)=a(d)t^2+2b_P(d)t+c_P,
\end{equation}
其中
\begin{equation}\label{eq:quadratic-radial-matrix-coefficients}
 a(d)=d^THd,\qquad
 b_P(d)=d^T\left(HP+\frac L2\right),\qquad
 c_P=\Phi(P).
\end{equation}
二次项只依赖方向和曲线的二次部分，一次项同时依赖基点与方向，常数项只依赖基点。改变极点，二次部分\(H\)不变；改变方向，基点处的值\(c_P\)不变。

由4.10.2的双线性表达，点\(P\)的极线在点\(X=P+td\)处满足
\begin{equation}\label{eq:quadratic-radial-matrix-polar}
 T(P,P+td)=c_P+t\,b_P(d)=0.
\end{equation}
这就是径向方程中的一次项与极线的联系。对\(c_P\ne0\)及非零交点参数，令\(z=1/t\)，则交点方程成为
\begin{equation}\label{eq:quadratic-reciprocal-radial-equation}
 c_Pz^2+2b_P(d)z+a(d)=0.
\end{equation}
两个倒数根的平均为\(-b_P(d)/c_P\)，而极线交点\(h\)满足\(1/h=-b_P(d)/c_P\)。式\eqref{eq:quadratic-radial-harmonic-mean}因而由二次型再次得到；4.11中的调和关系也正是在这条直线上对同一关系的射影表述。

\paragraph{1°中心与方向} \mbox{}

若\(P\)是中心，则\(HP+L/2=0\)，所有方向的一次项都消失。以中心为原点时，
\begin{equation}\label{eq:quadratic-centered-radial-form}
 t^2=-\frac{\Phi(P)}{d^THd},
\end{equation}
只要右边为正，就有关于中心对称的两个实交点。

对椭圆，化为标准方程后\(H=\operatorname{diag}(a^{-2},b^{-2})\)，方向向量为\((\cos\theta,\sin\theta)^T\)，所以
\begin{equation*}
 \frac1{r(\theta)^2}=\frac{\cos^2\theta}{a^2}+\frac{\sin^2\theta}{b^2}.
\end{equation*}
对双曲线，相应变为\(\cos^2\theta/a^2-\sin^2\theta/b^2\)。它大于零的方向有实半径，等于零的方向就是渐近线方向。一般中心双曲线的渐近线也由\(d^THd=0\)确定；在非中心的基点处，同一方向使径向二次项消失，但所得一次方程通常仍给出一个有限交点。

若\(d,e\)是一组互相垂直的单位向量，二次型在这组正交基下的两个对角系数为\(d^THd,e^THe\)。将\(d=(\cos\theta,\sin\theta)^T\)、\(e=(-\sin\theta,\cos\theta)^T\)代入，混合项相消，得到
\begin{equation}\label{eq:quadratic-orthogonal-radial-trace}
 d^THd+e^THe=H_{11}+H_{22}.
\end{equation}
因此，椭圆中两条互相垂直的中心半径的平方倒数之和为定值；普通部分中的三角恒等式，在这里表现为同一个二次型在不同正交基下的关系。

\par\addvspace{6pt}\noindent\begin{minipage}{\linewidth}
\begin{example}
以曲线\(5x^2+4xy+2y^2=1\)的中心为极点，写出\(r(\theta)\)，求半径长的最大值与最小值及相应方向。再证明任意两条互相垂直的中心半径满足\(r_1^{-2}+r_2^{-2}=7\)，将结果与旋转化简相比较。
\end{example}
\end{minipage}\par\addvspace{6pt}

\paragraph{2°焦点建系的二次型解释} \mbox{}

以焦点为原点、朝向准线的方向为\(x\)轴，式\eqref{eq:quadratic-focal-polar-cartesian}对应
\begin{equation*}
 H=\begin{pmatrix}1-e^2&0\\0&1\end{pmatrix},\qquad
 L=\begin{pmatrix}2e\ell\\0\end{pmatrix},\qquad
 F=-\ell^2.
\end{equation*}
所以其沿单位方向\(\widehat u=(\cos\theta,\sin\theta)^T\)的方程为
\begin{equation}\label{eq:quadratic-focal-radial-factorization}
 (1-e^2\cos^2\theta)t^2+2e\ell\cos\theta\,t-\ell^2
 =\bigl((1+e\cos\theta)t-\ell\bigr)
  \bigl((1-e\cos\theta)t+\ell\bigr)=0.
\end{equation}
两个有向根是\(\ell/(1+e\cos\theta)\)与\(-\ell/(1-e\cos\theta)\)。它们异号时交点在焦点两侧，同号时在同一侧；二次项为零时，只剩一个有限交点。这统一了前面双曲线焦弦中需要区分的方向。

原点的极线为\(e\ell x-\ell^2=0\)，即\(x=\ell/e=d\)，恰是准线。二次型的方向系数、焦点的极线及焦准线定义，由此共同给出同一个径向方程。

\paragraph{3°变换下的参数与长度} \mbox{}

设仿射变换为\(y=Kx+q\)，\(K\)可逆。射线\(P+td\)变为\(KP+q+tKd\)。参数\(t\)仍保留了线上点的线性关系，但\(Kd\)通常不是单位向量。若以像平面的实际有向距离\(s\)作参数，则
\begin{equation*}
 s=t|Kd|,\qquad \widehat d=\frac{Kd}{|Kd|}.
\end{equation*}
因此同一方向上各距离乘以同一个因子，倒数关系\(1/t_1+1/t_2=2/h\)在像图形中仍成立；不同方向的长度因子则通常不同。仿射变换保持割线、极线之间的关系，却不能据此把一个焦点直接当作像曲线的焦点。这与4.6中处理长度和垂直条件的原则相同。

\begin{exercise}
\mbox{}\par\nobreak
\begin{samepage}
（1）对\(x^2+2xy+3y^2-4x+2y-1=0\)，以\(P=(1,0)\)为极点。计算\(a(d),b_P(d),c_P\)，写出径向方程与\(P\)的极线。沿\(d=(1,0)^T\)核对交点的有向倒数关系。
\par\end{samepage}
\begin{samepage}
（2）设\(\Phi(P)\ne0\)。证明径向方程的判别式为\(4(b_P(d)^2-a(d)c_P)\)。若某方向有二重实根，说明对应直线为何为切线，并证明切点在\(P\)的极线上。椭圆\(x^2+2y^2=4\)内有一点\(P=(0,1)\)。过\(P\)的直线交椭圆于\(A,B\)，求\(|PA||PB|\)与弦长\(|AB|\)各自的取值范围，并指出最长弦的方向。
\par\end{samepage}
\begin{samepage}
（3）设椭圆\(x^THx=1\)，其中\(H\)正定。在任一组互相垂直的单位方向上取中心半径\(r_1,r_2\)，证明其平方倒数之和等于\(H_{11}+H_{22}\)。半径乘积\(r_1r_2\)是否也与方向无关？
\par\end{samepage}
\begin{samepage}
（4）将\(y^2=4x\)平移到以焦点为原点的坐标系，写出\(H,L,F\)，再由式\eqref{eq:quadratic-radial-matrix-equation}分解径向方程，得到焦点极坐标表示。
\par\end{samepage}
\begin{samepage}
（5）在非退化中心双曲线中，设基点\(P\)不在曲线上。对渐近线方向\(d\)，讨论\(b_P(d)\ne0\)与\(b_P(d)=0\)时的有限交点个数，并解释后一情形中\(P\)与相应渐近线的位置关系。
\par\end{samepage}
\begin{samepage}
（6）伸缩\((x,y)\mapsto(2x,y)\)将单位圆变为椭圆。取圆内非中心点\(P\)，过\(P\)作一条与其极线相交的割线，证明交点与极线交点的有向倒数关系在伸缩后保持。再比较圆心的像与像椭圆的两个焦点，说明焦点不能直接通过这个伸缩得到。
\par\end{samepage}
\end{exercise}
